\uuid{U2iA}
\titre{Distance}
\auteur{bodin ; hanani ; mzari}
\organisation{exo7}
\licence{CC BY-NC-SA 4.0}
\datecreate{2020-01-20}
\niveau{L1}
\module{Géométrie}
\chapitre{Géométrie du plan}
\sousChapitre{Distance, norme, produit scalaire}
\difficulte{2}
\theme{distance}
\exo7qcmid{lille-1.geometrie-plane.019}

\contenu{
\begin{question}
\qclassification{140.01,140.02}
\qtags{motcle=distance}
Dans le plan rapporté à un repère orthonormé, on considère la droite $D$ d\'efinie par le paramétrage :
$$\left\{\begin{array}{ccl}x&=&1+t\\ y&=&2-t,\quad t\in \Rr.
\end{array}\right.$$
Quelle est la distance $d$ du point $M(2,3)$ à la droite $D$ ?
\begin{answers}
    \good{$d=\sqrt{2}$}
    \bad{$d=\sqrt{3}$}
    \bad{$d=1$}
    \bad{$d=2$}
\end{answers}
\begin{explanations}
Le point $A(1,2)\in D$ et le vecteur $\vec{v}=(1,-1)$ dirige $D$. On utilise la formule $\displaystyle d=\frac{\left|\mbox{det}\left(\overrightarrow{AM},\vec{v}\right)\right|}{\|\vec{v}\|}=\sqrt{2}$.
\end{explanations}
\end{question}
}
